% NAF 5176, batch 1 — Gallica views 1–20.
% Pass: Opus 5.5 (claude-opus-5-5), 2026-09-22 — first pass, unchecked against the views by a human.
%
% What the twenty views hold. Colour scans (converted to grey for reading),
% 8734 × 6362 for the openings, 4286 × 6509 for the front board. From view 2
% on, every view is the volume lying open: a left page and a right page, the
% gutter near x = 4350–4700. Nine views carry text in the writers' hands.
%
%   v. 1       Front board, marbled paper. Not transcribed.
%   v. 2       Contreplat with the printed oval label « FR / NOUV. ACQ. /
%              5176 »; blank flyleaf. Not transcribed.
%   v. 3–5     Blank flyleaves and guards (checked at full contrast).
%   v. 6       A modern slip on squared paper pasted to a leaf, dated
%              « 21 mai 1962 », signed with a paraph not read: a reader's
%              note that some text of the volume corresponds to Huygens's
%              « de ratiociniis in ludo aleae », probably a French
%              translation earlier than the version given to Van Schooten,
%              « Date possible : 1656 », « Copiste et traducteur possible :
%              Des Billettes ». Not transcribed (a twentieth-century note);
%              the text it describes is not in this batch.
%   v. 7       The library's title leaf, 1888: « Papiers de Mersenne & de
%              Pascal », « Volume de 42 Feuillets plus le Feuillet A
%              préliminaire. Les Feuillet 22, 24, 29, 35, 36 sont blancs.
%              29 Septembre 1888. », « Libri 1848. », « Fr. nouv. acq.
%              5176. ». Not transcribed.
%   v. 8       The « Feuillet A »: an ink « A » top right, and a
%              nineteenth-century note « Manuscrits autographes du père
%              Mersenne », with « Libri 1848 » and « FR. nouv. acq. 5176 ».
%              Not transcribed.
%   v. 9       f. 1r. « Preface. » of a French treatise on composite
%              motions (mouvemens composez, impetuosité, reflexion), in a
%              very fast, heavily corrected hand; « 1 Proposition ».
%   v. 10      The same opening as v. 9, photographed a second time. Not
%              transcribed.
%   v. 11      f. 1v (« II Proposition », equality of composite motions
%              on a circle and on its tangent) | f. 2r (« Lemme » on
%              uniformly decreasing speed; « Proposition III », the first
%              turn of Archimedes's spiral equals the line whose square is
%              the sum of the squares of the circumference and the radius).
%   v. 12      f. 2v (end of prop. III, the rest of the page blank) |
%              f. 3r, a second, smaller and denser hand: a table of chapters
%              of a treatise « Des triangles Rectilignes et Spheriques »
%              (ch. 1–25, with page references), then problems numbered
%              1, 2 on ellipses whose axes and focal distance are rational,
%              from Pythagorean triangles (44.117.125, 35.120.125,
%              75.100.125; 16.63.65, 25.60.65, 33.56.65, 39.52.65).
%   v. 13      f. 3v (problems 3–6 of the same series: 3.4.5, 15.20.25,
%              8.15.17, 125² = 15625, 65² = 4225, 8450) | f. 4r (problems
%              7–8: the lines of an ellipse from the « triangle
%              fondamental » 8.15.17, a column of eight products, and eleven
%              lines of an ellipse made integral by multiplying by 7:
%              AC 490 … LE 175). The same hand as f. 3.
%   v. 14      f. 4v blank (show-through, not transcribed) | f. 5r, the
%              fast hand of f. 1–2 again: the end of a passage on the
%              earthshine on the Moon, then « Proposition III », « encor
%              autrement », why reflection is at equal angles, with a
%              rigid physical line AB striking a plane.
%   v. 15      f. 5v, the same argument continued, a wide ink blot across
%              the middle of the page; « Avertissement pour la Reflexion
%              des rayons ». On a guard at the right, two small printed
%              figures, cut out and pasted, numbered « 6 » and « 7 » (a cut
%              cone; the conic sections round a common vertex A).
%   v. 16      f. 8r (the figures 6 and 7 folded back over its margin):
%              « Proposition » (numbered 8?), « Expliquer toutes les
%              Sections Coniques »; the writer refers to his own
%              treatments of the conics « depuis la page 498 du
%              Commentaire sur la Genese iusques à la 535 », « le 16
%              chapitre du 4 livre de la verité des sciences », « le
%              premier livre de l'Harmonie universelle, prop. 28 », « la
%              5 propos. de l'Utilité de l'Harmonie », « la 18 et 19
%              propos. des Hydrauliques ». These are the titles of Mersenne's
%              books (Quaestiones in Genesim, La Vérité des sciences,
%              Harmonie universelle, the Hydraulica of the Cogitata); the
%              leaf itself names no author. That the fast hand is
%              Mersenne's is an inference from these first-person
%              references, and from the 19th-century note on the leaf A,
%              not a reading.
%   v. 17      f. 8v with three printed figures pasted in its margin (« 9 »
%              parabola, « 10 » ellipse, « 11 » hyperbola) covering the
%              beginnings of the lines, and the narrow slip f. 12 unfolded
%              at the right. Described in a note, not transcribed.
%   v. 18      The same opening with the figures lifted: f. 8v,
%              « Proposition. Expliquer des manieres bien aisées pour
%              descrire les trois sections coniques » (by foci and
%              compass; Mydorge's Sections coniques cited), the sentence at
%              its foot continued on the slip f. 12r (five lines).
%   v. 19      Verso of the slip f. 12 (blank) and a grey guard leaf
%              (blank, stamps showing through). Not transcribed.
%   v. 20      Grey guard (blank) | f. 15r, the fast hand, with printed
%              figures « 13 » and « 14 » pasted on it: « 1 Proposition »
%              on refraction (given one refracted ray, find all the
%              others), « Corollaire », « II Proposition », which runs off
%              the foot of the page in mid-sentence (« plus rare, et plus
%              subtil que »).
%
% Who writes to whom. No letter in this batch: no salutation, address,
% signature or date anywhere in views 9–20. The only dates in the views are
% the library's (« 29 Septembre 1888 », « Libri 1848 ») and the modern
% note's (« 21 mai 1962 »).
%
% Two hands. (1) A very fast French cursive, with long connecting strokes,
% very many deletions and interlinear additions, words often left without
% lifting the pen (f. 1–2, 5, 8, 12, 15). Many of its words are left
% \ill{} or \uncertain{}; the densest passages are f. 5r lines 1–6, f. 5v
% (under the blot) and the paragraph on the glass cylinders there. (2) A
% smaller, closer and more abbreviated hand, with many figures, on f. 3–4
% only (the triangle chapters and the rational ellipses); its numbers read
% well and its words much less well.
%
% The printed figures. Small woodcut or engraved diagrams, each cut out,
% pasted on a leaf or a guard, and numbered in ink 6, 7, 9, 10, 11, 13,
% 14 — numbers of the same series as the leaves' foliation. They are not
% drawn in the manuscript; they are described in \note{}, with their
% lettering, and not reproduced. They look like proofs or offprints of the
% plates of a printed book on optics and the conics; which book is not
% established here.
%
% Notation. None beyond letters of figures, integers and the fractions
% 17 1/7, 52 6/7 (set $\frac17$, $\frac67$). « R » or « R. » is a barred R,
% the sign of the root, kept as R. « $\square$ » is a small square for « quarré »
% (f. 3v), kept as $\square$. « q. » abbreviates « quarré ». The writers' « | »
% between results is kept. A long s is transcribed as s. The fast hand's
% « & » and « et » are both « et ».
%
% Numbers on the leaves.
% - Foliation. A bold ink figure at the top right of each recto, one
%   series across the volume: 1 (v. 9), 2 (v. 11 R), 3 (v. 12 R), 4
%   (v. 13 R), 5 (v. 14 R), 8 (v. 16 R), 12 (the slip, v. 17–18), 15
%   (v. 20 R, underlined). The missing figures are carried by the pasted
%   printed figures (6, 7 at v. 15; 9, 10, 11 at v. 17; 13, 14 at v. 20),
%   which confirms a single count of every piece of paper, and it agrees
%   with the 1888 title leaf (« Volume de 42 Feuillets plus le Feuillet A
%   préliminaire »), whose « A » is written the same way on v. 8. This is
%   the series the library counted in 1888. It is ink, not pencil, so, as for
%   Latin 17859 and Français 9115, no \folio{} is written until a human has
%   looked at the leaves (coordinator's reconciliation of the three
%   batches; see batch 2's header for the whole series and the certificate's
%   list of blank leaves, which it matches). No verso carries a number.
% - An older, finer ink series beside it on the same rectos, not in
%   order: 41 (f. 1), 42 (f. 2), 10 (f. 3), 11 (f. 4), 19 (f. 5), 14
%   (f. 8), 15 (f. 12), 16 (f. 15). Probably the numbering of an earlier
%   arrangement (Libri's?); an inference. Recorded in notes only.
% - Small figures at the top middle of the pages of the treatise on
%   motions: 5 (? v. 9, beside « Preface. »), 6. (f. 1v), 7. (f. 2r), 8.
%   (f. 2v) — the writer's own pagination, apparently. Recorded in notes.
% - « 20 » in ink on the edge of a paper projecting above f. 15 at v. 20:
%   not this leaf's number.
% - View 20 is re-shot at view 21 with its pasted figures lifted; its text is
%   given whole there, in batch 2, and only described here.
% - Marginal references in large characters, boxed: « 2. M. fol. 86 »
%   (f. 5r), « 3. M. fol. 87 » (f. 5v). Transcribed as \marginal{}; their
%   meaning (a volume « M » and its folios?) is not established.
% - Problem numbers 1–8 in the left margin of f. 3–4, and the writer's
%   proposition numbers, are content.
% - Shelfmarks and stamps: « FR. nouv. acq. 5176 », « R.c. 8070 (88) »,
%   the oval « ACQ. 8070 (LIBRI) » (f. 1), round « BIBLIOTHÈQUE NATIONALE
%   MSS / R.F. » stamps throughout.
%
% The catalogue's dating is copied verbatim.

\documentclass[11pt,a4paper]{article}
% The preamble shared by every transcription of the mathematicians' manuscripts in Gallica.
%
% It rests on one idea: **uncertainty compounds, it does not resolve.** A
% transcription that smooths over a doubtful word has destroyed the only thing
% it was made for — a reader must be able to tell, at every point, what was on
% the page and what was supplied. The apparatus macros below are therefore the
% critical apparatus, not decoration.
%
% The same commands are recognised by `scripts/render.mjs`, which produces the
% site's reading view, and by `scripts/tei.mjs`. A macro added here must be
% added there too, or rendering fails loudly — which is the wanted behaviour.
%
% Comments here are English, like the rest of the repository. The documents
% this preamble sets are French, and that is deliberate: see the skills.

\ProvidesPackage{germain}[2026/09/15 Transcriptions of mathematicians' manuscripts in Gallica]

\usepackage{amsmath,amssymb,amsthm}
\usepackage{mathrsfs}
% Kept from the parent preamble so the renderer's diagram support stays
% available; these manuscripts draw no commutative diagrams, and nothing here uses it.
\usepackage{tikz-cd}
\usepackage[english,french]{babel}
\usepackage{fontspec}

% --- Greek ------------------------------------------------------------------
% The text face stays fontspec's default, Latin Modern, which has no Greek:
% XeTeX drops every glyph it lacks with a line in the log and nothing on the
% page, and the Greek words the manuscripts quote vanished from the PDFs. So
% the Greek and Greek Extended blocks get a XeTeX character class of their own,
% and entering or leaving a run of it switches the family to CMU Serif — the
% same Computer Modern design, with polytonic Greek — and back. Only the
% family changes: a Greek word in \textit or \textbf keeps its shape and series.
% The transitions to and from every other class carry the tokens that class
% has towards an ordinary letter, so French spacing before « ; » or inside
% « » still holds after a Greek word. No grouping, so no brace or \endgroup
% can come out unbalanced; maths is left alone, since XeTeX inserts nothing
% there. Kept identical in `exercices.sty`.
\makeatletter
\newfontfamily\germainGreekfont{cmun}[
  Extension=.otf, NFSSFamily=germaingreek,
  UprightFont=*rm, ItalicFont=*ti, SlantedFont=*sl,
  BoldFont=*bx, BoldItalicFont=*bi, BoldSlantedFont=*bl]
\newXeTeXintercharclass\germain@greek
\def\germain@greekrange#1#2{%
  \@tempcnta=#1\relax
  \loop
    \XeTeXcharclass\@tempcnta=\germain@greek
    \ifnum\@tempcnta<#2\advance\@tempcnta\@ne\repeat}
\germain@greekrange{"0370}{"03FF}
\germain@greekrange{"1F00}{"1FFF}
\def\germain@greekprev{\rmdefault}
\def\germain@greekin{%
  \edef\germain@greekprev{\f@family}\fontfamily{germaingreek}\selectfont}
\def\germain@greekout{\fontfamily{\germain@greekprev}\selectfont}
\def\germain@greekjoin{%
  \XeTeXinterchartokenstate=\@ne
  \@tempcnta=\z@
  \loop
    \ifnum\@tempcnta=\germain@greek\else
      \XeTeXinterchartoks\germain@greek\@tempcnta=\expandafter{%
        \expandafter\germain@greekout\the\XeTeXinterchartoks\z@\@tempcnta}%
      \XeTeXinterchartoks\@tempcnta\germain@greek=\expandafter{%
        \the\XeTeXinterchartoks\@tempcnta\z@\germain@greekin}%
    \fi
    \ifnum\@tempcnta<4095\advance\@tempcnta\@ne\repeat}
% babel-french sets its own classes, and saves and restores the interchar
% state, each time a language is selected: join again after it does.
\addto\extrasfrench{\germain@greekjoin}
\addto\extrasenglish{\germain@greekjoin}
\AtBeginDocument{\germain@greekjoin}
\makeatother

\usepackage{geometry}
\geometry{a4paper,margin=25mm,marginparwidth=28mm}
\usepackage{marginnote}
\usepackage{xcolor}
\usepackage[normalem]{ulem}
\setlength{\emergencystretch}{3em}
\usepackage{hyperref}
\hypersetup{colorlinks=true,linkcolor=black,urlcolor=black,pdfborder={0 0 0}}

% --- Batch metadata ---------------------------------------------------------
% Taken as they stand by the rendering script, which shows them at the head of
% the reading view. They fix what the file claims to cover. The macro names are
% the parent project's, so that the scripts stay shared; \folder here means the
% volume's slug — `fr-9115` — and \pages counts Gallica views.

\newcommand{\folder}[1]{\gdef\grothFolder{#1}}
\newcommand{\batch}[1]{\gdef\grothBatch{#1}}
\newcommand{\pages}[2]{\gdef\grothFirst{#1}\gdef\grothLast{#2}}
\newcommand{\foldertitle}[1]{\gdef\grothFoldertitle{#1}}

% The holder's own shelfmark, printed as they write it: « Français 9115 ».
\newcommand{\shelfmark}[1]{\gdef\grothShelfmark{#1}}

% The dating, copied verbatim from the holder's catalogue — never inferred
% here. « XIXe siècle » is the BnF's; a pass that dates a leaf from its content
% says so in a \note{}, as its own inference.
\newcommand{\dating}[1]{\gdef\grothDating{#1}}

% The status notice, printed in the title block and repeated as a diagonal
% watermark on every page of the PDF. The manuscripts are in the public
% domain; what this line declares is not a right but a quality — an
% unverified first pass by a machine. The skills write the line; a file
% without it gets no watermark, so leaving it out is visible in the output.
\newcommand{\watermark}[1]{\gdef\grothWatermark{#1}}

\gdef\grothFolder{?}\gdef\grothBatch{}\gdef\grothFirst{?}\gdef\grothLast{?}%
\gdef\grothFoldertitle{}\gdef\grothShelfmark{}\gdef\grothDating{}\gdef\grothWatermark{}

% --- The résumé -------------------------------------------------------------
\newenvironment{resume}
  {\small\setlength{\parskip}{0.45em}}
  {\par\medskip\hrule\medskip}

% --- Keywords ---------------------------------------------------------------
% In English, closing the modernised reading's résumé: the modern vocabulary
% under which to search for what the leaves construct. The manifest extracts
% them to tag the volume — this line is the single source of the tags.
\newcommand{\keywords}[1]{\par\smallskip\noindent
  {\footnotesize\textsc{Keywords} --- \textit{#1}}\par}

% --- The critical apparatus -------------------------------------------------

% The Gallica view number, in the margin. This is the anchor that lets a
% disagreement over a formula be settled by looking at *one* image rather than
% a volume of 750. The site uses it to turn the facsimile as the transcription
% is scrolled. A view is one scanned image: usually one side of a leaf.
\newcommand{\page}[1]{%
  \marginnote{\footnotesize\color{gray!70}\textsf{v.\,#1}}%
  \label{page:#1}\ignorespaces}

% The BnF's pencilled foliation, where the leaf carries it and a pass can read
% it: « 198r ». This is what the literature cites — « ff. 198r–208v » — and
% the one line that turns a citation into a view. Recorded, never inferred:
% a folio a pass cannot read is not written.
\newcommand{\folio}[1]{%
  \marginnote{\footnotesize\color{gray!50}\textsf{f.\,#1}}\ignorespaces}

% The modernised reading's coarser counterpart to \page{}: which views a
% section is drawn from.
\newcommand{\pagerange}[2]{%
  \marginnote{\footnotesize\color{gray!70}\textsf{v.\,#1--#2}}%
  \ignorespaces}

% Illegible: never guessed.
\newcommand{\ill}{\textcolor{red!60!black}{[\,\ldots\,]}}

% A doubtful reading: the word is given, and flagged as uncertain.
\newcommand{\uncertain}[1]{\uline{#1}}

% An editorial addition — a missing letter, an implied word.
\newcommand{\add}[1]{\textcolor{blue!50!black}{[#1]}}

% Struck out by the writer. What was crossed out is part of the document.
\newcommand{\struck}[1]{\sout{#1}}

% The transcriber's note, on the state of the page or the mathematics.
\newcommand{\note}[1]{\par\smallskip{\small\color{gray!60!black}\textit{[#1]}}\par\smallskip}

% A marginal note of hers, distinct from the body of the text.
\newcommand{\marginal}[1]{\par\smallskip\noindent\hspace{1em}%
  {\small\color{gray!60!black}#1}\par\smallskip}

% --- Title ------------------------------------------------------------------
% The title block is French, like the documents themselves.

\AddToHook{shipout/background}{%
  \ifx\grothWatermark\empty\else
    \begin{tikzpicture}[remember picture, overlay]
      \node[rotate=52, scale=3.4, align=center, text=gray!18,
            font=\sffamily\bfseries]
        at (current page.center) {\grothWatermark};
    \end{tikzpicture}%
  \fi}

\AtBeginDocument{%
  \begin{center}
    {\large\bfseries Manuscrits de mathématiciens (Gallica) ---
      \ifx\grothShelfmark\empty\grothFolder\else\grothShelfmark\fi}\\[2pt]
    {\itshape\grothFoldertitle}\\[4pt]
    {\small vues \grothFirst--\grothLast
      \ifx\grothBatch\empty\else\ (lot \grothBatch)\fi}\\[2pt]
    \ifx\grothDating\empty\else
      {\small\color{gray!60!black}Datation du catalogue : \grothDating}\\[2pt]
    \fi
    \ifx\grothWatermark\empty\else
      {\small\bfseries\color{red!45!gray}\grothWatermark}\\[2pt]
    \fi
    {\footnotesize\color{gray!60!black}Transcription établie d'après les images IIIF de Gallica.
     Source gallica.bnf.fr / Bibliothèque nationale de France.\\
     Les lectures incertaines sont soulignées, les illisibles notées [\,\ldots\,],
     les ajouts entre crochets ; \textsf{v.} numérote les vues de Gallica,
     \textsf{f.} la foliotation au crayon de la BnF.}
  \end{center}
  \vspace{1em}}

\endinput


\folder{naf-5176}
\shelfmark{NAF 5176}
\batch{1}
\pages{1}{20}
\dating{1601-1700}
\watermark{Lecture automatique\\non vérifiée}
\foldertitle{Papiers du Père MERSENNE et de Blaise PASCAL.}

\begin{document}

\page{9}
\note{Premier feuillet du texte, monté sur onglet. En haut, au milieu, « Preface. », de la main du texte ; à sa droite un petit signe qui se lit 5. En haut à droite, à l'encre forte, le chiffre « 1 » (série continue ; voir l'en-tête), et plus à droite, d'un trait plus fin, « 41 », d'une numérotation plus ancienne ; au-dessous, dans une accolade, la cote « FR. nouv. acq. 5176 ». Au pied, dans une accolade, « R.c. 8070 (88) », et le cachet ovale « ACQ. 8070 (LIBRI) » ; au milieu de la page, le cachet rond « BIBLIOTHÈQUE NATIONALE MSS » avec « R.F. ». La vue 10 reproduit la même ouverture, photographiée une seconde fois : elle n'est pas transcrite à part.}

Preface.

Entre plusieurs choses qui servent à relever la puissance de la lumiere, et de ses rayons il faut donner le premier rang aux mouvemens tant simples que composez, dont l'usage est si general estant ce qui \uncertain{anime} les mechaniques et \ill{} \ill{}, qu'il y a bien peu de choses où \uncertain{elles} \ill{} par \ill{} : c'est pourquoy ie \uncertain{desire} preparer le lecteur à cette sorte de \struck{\ill{}} speculation, qui peut aussi bien servir aux Theologiens qu'aux Philosophes, \uncertain{que} et aux orateurs qui voudront s'en servir.
\note{écriture très rapide, à longues liaisons ; la ligne 3 est la plus douteuse. Après « les mechaniques », deux mots commençant par « et le » n'ont pu être lus ; la fin de la ligne 3, au bord noirci du feuillet, et le mot qui ouvre la ligne 4 après « par » restent illisibles. « qu'aux Philosophes » est suivi d'un mot bref à longue queue de q.}

Soit par exemple le mouvement AC composé du mouvement AD, et AB, et qu'on veuille scavoir \struck{quelle} \uncertain{raison} quel rapport, ou quelle raison il a avec les \add{deux} mouvemens qui le composent. Je di que sa vitesse est à celle d'AD, comme la ligne AC à la ligne AD ; et à celle d'AB, comme la ligne AC à la ligne AB, puisqu'ils se font en \uncertain{mesme} temps.
\note{« deux » est écrit au-dessus de la ligne, avec un signe d'insertion. En marge de gauche, à hauteur de ce paragraphe et du suivant, deux figures : un rectangle ABCD (lettres A, B, C, D, E, F, G, H) avec la diagonale AC, les perpendiculaires et un segment pointillé ; au-dessous, un demi-cercle chargé de hachures et de traits repassés, lettré B, H, G, M, D, E, F, C, dont une partie est illisible.}

Quant à l'impetuosité, \add{elle} est autre chose que la vitesse considerée en tel point \struck{de la ligne} qu'on voudra, puisque dans la ligne \struck{du mouvement} descrite par le mouvement : de sorte qu'elle est tousjours egale dans \struck{\ill{}} le mouvement \ill{},
\note{sous « la ligne \ldots{} descrite », une ligne interlinéaire en petite écriture, qui paraît se terminer par « de l'impetuosité », n'a pu être lue en entier. La phrase paraît incomplète ; le dernier mot de la ligne, qui qualifie « le mouvement », n'a pas été lu.}

Dans cette figure l'impetuosité \uncertain{receüe} au point C du mouvement qui vient d'A par le mouvement composé, est à l'impetuosité \uncertain{receüe} en D du mouvement \add{qui vient} \struck{\ill{}} d'A, comme la ligne AC à la ligne AD. Il faut dire la mesme chose de l'impetuosité C comparée à l'impetuosité B.

Et si FE est perpendiculaire à AC, l'effort qui produira l'impetuosité \struck{\ill{}} \add{composée} prise en C sera à l'effort qui produira l'impetuosité composée prise en D ou B, comme AC à AD, ou AB. Lors qu'on voudra scavoir quelle partie \struck{\ill{}} du mouvement AC respond à ce qu'il vient de DA, et de AB, il faut \struck{\ill{}} tirer une droite de \struck{\ill{}} l'angle B sur AC, qui divise l'angle ABC en deux parties egales, comme fait la ligne BD, car l'angle \struck{EBF} \add{CBD} est egal à l'angle \struck{FBG} \add{DBA} : de \struck{\ill{}} maniere que la ligne qui de B tombe perpendiculairement sur la ligne AC \add{au point H,} sert \struck{\ill{}} \add{\ill{}} pour cela, comme i'ay remarqué dans la 32 prop. de la Ballistique.
\note{« CBD » et « DBA » sont écrits au-dessus de « EBF » et « FBG » biffés. Au-dessus du mot biffé qui suit « sert », un mot bref n'a pas été lu. « la Ballistique » renvoie, selon toute apparence, à un ouvrage de l'auteur ; la pièce ne le nomme pas autrement.}

Mais ie veux aiouter quelque proposition \struck{des} \add{de} l'excellent esprit qui fonde sur cette hypothese que deux mouvemens s'opposent, et se retardent d'autant plus qu'ils sont plus opposez, \struck{et que cette opposition} \struck{\ill{}} que cette opposition est plus grande dans le \struck{\ill{}}
\note{« l'excellent esprit » est lu avec doute : le mot a été récrit après la rature de « des ». Un long trait court sous les deux dernières lignes de ce paragraphe ; il est lu ici comme un soulignement, les ratures étant les traits plus lourds qui barrent « et que cette opposition », le premier mot de la ligne suivante et le dernier mot, pris dans le cachet de la bibliothèque. La phrase reste inachevée.}

1 Proposition

Quand \struck{\ill{}} deux mouvemens droits uniformes faits sur deux lignes qui se rencontrent, ils \add{s'empeschent et} se retardent d'autant plus que l'angle qu'ils font est plus grand, et moins si il est plus petit. La mesme chose arrive si \struck{le mouvement} \struck{se} la ligne du mouvement droit rencontre la circonference d'un cercle.
\note{l'énoncé est enfermé à gauche dans une grande accolade. « s'empeschent et » est écrit au-dessus de « se retardent ».}

Soit la droite \add{perpendiculaire} AC, \struck{\ill{} perpendiculaire \ill{}}, sur laquelle se fera le mouvement d'A \struck{\ill{}} \add{vers} B ; \uncertain{et puisque} les droites DB, EB, FB tombent au point B : que l'angle \uncertain{ADB} soit moindre que l'angle ABE, et ABE moindre que l'angle ABF, et que le mouvement qui se fait sur DB, EB, FB soit \uncertain{egal}, ie di que le mouvement sur AB et sur DB se \uncertain{nuisent} et retardent moins, que les deux mouvemens sur AB et sur EB, et ceux-cy \struck{se \uncertain{nuisent} moins} s'empeschent moins, que les deux sur AB, FB. : pourceque l'angle \uncertain{ADB} est moindre qu'\uncertain{AED}, dont les mouvemens sont moins opposez, et partant ils sont moins retardez.
\note{en marge de gauche, la figure : une verticale AC coupée en B, un arc de cercle ABC, et les droites DB, EB, FB aboutissant en B, D relié à G par un pointillé. Les noms d'angles « ADB » et « AED » sont lus tels qu'ils paraissent ; la figure et l'énoncé demanderaient ABD et ABE.}

De là il s'ensuit que le mouvement composé des mouvemens sur AB et FB est plus tardif que \struck{celuy qui se} le composé du mouvement sur AB et EB, et que celuy-cy est plus tardif que \struck{celuy qui se fait} \add{le} composé des mouvemens sur AB, DB.

\page{11}
\note{La vue montre une ouverture : à gauche le verso du feuillet de la vue 9, à droite le recto du feuillet suivant (donné plus bas). En haut de la page de gauche, au milieu, « 6. » ; aucun autre chiffre.}

II Proposition.

Si le mobile \struck{\ill{}} qui est meu par la circonference d'un cercle d'un mouvement uniforme de quelque vitesse que ce soit, et que pendant qu'il se meut il luy survienne un autre mouvement uniforme de \add{telle} quelque vitesse \struck{qu'il voudra par \ill{}} et \add{par} tel angle qu'il voudra, et que de ces deux mouvemens il s'en fasse un mouvement composé. Et de mesme, si le mobile se meut par une ligne droite de la \uncertain{pareille} vitesse dont il se mouvoit sur la circonference du cercle, et que pendant qu'il se meut, il luy survienne un autre mouvement droit en \uncertain{semblable} angle, et de la mesme vitesse qu'avoit le mouvement droit qui estoit survenu au mouvement circulaire, et que de ces deux mouvemens il s'en compose un autre, les vitesses de ces deux mouvemens composez seront egales entr'elles.
\note{l'énoncé est enfermé à gauche dans une grande accolade. « telle » et « par » sont écrits au-dessus de la ligne.}

\struck{Soit} Du centre B soit descrite telle circonference de cercle qu'on voudra DCE, dans laquelle le mobile se meuve de D à C de telle vitesse uniforme qu'on voudra : soit \struck{\ill{}} \add{la droite} \struck{\ill{}} EC qui touche le cercle DC au point C. Et que AC soit tirée qui fasse tel angle qu'on voudra. Soient les angles ACE, ACD soient egaux, ou du moins ACE ne soit pas plus grand qu'ACD.
\note{en marge de gauche, la figure : des droites issues de C vers A, G, E, H, D, F, I, un arc de cercle passant par D, K, C, et un pointillé de C à B. Le mot qui suit « la droite », écrit au-dessus de la ligne, est lu avec doute.}

Que le mobile \struck{\ill{}} se meuve de E à C, de la pareille vitesse \ill{} comme nous avons supposé qu'il se meut de D à C : et qu'un autre mobile se meuve aussi d'A à C uniformement \struck{\ill{}} de telle vitesse qu'on voudra : il se fera un mouvement composé des mouvemens sur AC, et DC ; et encore un autre mouvement composé des deux qui se font sur AC et EC.

Or ie di que ces deux mouvemens \add{ainsi} composez sont egaux entr'eux en egale vitesse. Car s'ils ne sont pas egaux, le composé d'AC et DC sera plus ou moins grand que le composé d'AC et EC. Et qu'il soit donc plus grand, \struck{si faire se peut} s'il est possible. Que FC soit tirée faisant l'angle FCD plus grand que l'angle ACD, et lors (par la precedente) le composé sera moindre : faisons donc que le mouvement composé des mouvemens FC, DC ne soit pas plus grand, mais egal au composé des mouvemens AC, EC : donc, puisque l'angle FCE est plus grand que l'angle ACE, le mouvement composé des mouvemens FC, EC, est moindre que le composé d'AC, EC. Mais l'angle FCD est egal à l'angle FCE, \struck{\ill{}} \add{\ill{}} il est \uncertain{aussi pas moindre}, donc les mouvemens sur FC, DC \struck{sont plus \ill{} \ill{}} \struck{sont plus} \add{ne} sont \uncertain{pas} moins opposez que les mouvemens sur FC, EC.
\note{ce passage est chargé de ratures superposées, écrites au-dessus et au-dessous de la ligne ; l'ordre des mots retenu est une lecture, non une certitude.}

Donc le composé de FC, DC est egal ou moindre que le composé de FC, EC : mais ce composé a esté \add{prouvé} \struck{\ill{}} estre moindre que le composé d'AC, EC, donc le composé de FC, DC, est moindre que le composé d'AC, FC. Et neantmoins il estoit supposé egal, et partant il sera egal et inegal, ce qui est absurde. Donc le mouvement composé des mouvemens sur AC et DC, n'est pas plus grand que le composé \struck{d'AC des mouvemens} AC, EC.
\note{« d'AC, FC » : ainsi écrit ; l'argument demande AC, EC.}

De mesme ; que le composé des mouvemens sur AC, DC soit moindre (s'il est possible) que le composé d'AC, EC. Donc menons GC, afin que \struck{GCD soit \ill{}} l'angle GCD soit moindre que l'angle ACD, et partant que le mouvement composé soit plus grand, que le composé de GC, DC soit egal au composé d'AC, EC : soit tirée la secante CH qui coupe la circonference au point K, et que l'angle ECH soit moindre que l'angle \uncertain{ACE}, donc le composé de GC, HC sera plus grand que le composé d'AC, EC, pourceque l'angle GCH est moindre que l'angle ACE. Mais le composé des mouvemens sur GC, DC est plus grand que le composé des mouvemens sur GC, HC (pourceque l'angle GCD, \struck{\ill{}} qui est fait par la courbe KD et par la droite GC, est moindre que l'angle GCH fait \struck{par} par les droites GC et CK) donc le mouvement composé des mouvemens GC, DC est plus grand que le composé \struck{des GC, DC} des mouvemens AC, EC. Et neantmoins il estoit posé egal, et partant il sera egal et inegal, ce qui est impossible. Donc le mouvement composé des mouvemens sur AC, DC n'est ni moindre ni plus grand que le composé des mouvemens sur AC, EC, donc les vitesses de ces deux mouvemens composez sont egales.
\note{la dernière lettre de « l'angle \uncertain{ACE} » est couverte d'une tache d'encre.}

\note{Page de droite de la vue 11 : recto du feuillet suivant. En haut au milieu, au-dessus du premier paragraphe, « Lemme » et, à sa droite, « 7. » ; en haut à droite, à l'encre forte, « 2 » (série continue ; voir l'en-tête), et plus à droite, d'un trait plus fin, « 42 ».}

Lemme

La vitesse qui se diminue tousiours iusqu'à ce qu'elle cesse, et \ill{} que le temps dans lequel se fait le mouvement, \ill{}, est la moitié de la plus grande vitesse qui a esté au commencement du mouvement.
\note{énoncé enfermé à gauche dans une accolade. Le mot qui suit « cesse, et » (six ou sept lettres, avec un s long) et celui qui suit « le mouvement, » n'ont pu être lus.}

Que la vitesse proposée soit celle dont le mobile se meut uniformement sur la droite AB, dans le temps que le doigt \struck{\ill{} sur le \ill{}} parcourt la circonference CD : et que \struck{cette} \add{la} circonference DE soit egale à DC, afin que CE soit double de DE. Et pourceque le mobile qui se meut uniformement requiert le temps DC pour parcourir AB, quand sa vitesse se diminue, il requiert un \uncertain{moindre} temps, de sorte que \struck{le temps} le temps doit augmenter autant que la vitesse se diminue. Et par consequent si la \uncertain{parfaite} \struck{\ill{}} diminution se fait sur AB, le temps CD doit estre augmenté d'autant que la vitesse diminue ; donc le mobile \struck{\ill{}} ira sur AB \struck{\ill{}} \add{avec} quelque degré de vitesse dans tout le temps de CD et puis il se reposera. Et partant il employera tout ce temps qui est double du temps CD, qu'il employeroit seulement lorsqu'il conservoit toute sa premiere vitesse sur AB : d'où il faut conclure que cette vitesse demeurant entiere est double de \uncertain{la moyenne} qui va se diminuant.
\note{en marge de gauche, une petite figure : la droite AB, et à côté une ligne CDE coupée par un trait vertical en D. « le doigt » désigne, semble-t-il, l'aiguille d'un cadran qui mesure le temps. Au-dessus de « avec », un mot biffé commençant par « le ».}

Proposition III.

\struck{La spirale} La premiere revolution de la spirale d'Archimede est egale à la \struck{demi-circonference} droite qui est egale en puissance à la circonference et au semidiametre de son cercle.
\note{énoncé enfermé à gauche dans une accolade.}

Soit descrit du rayon AB tel cercle qu'on voudra BCDE : \struck{soient des rayons} soient tant de poincts qu'on voudra F, G, H \struck{dans les rayons}, qui descrivent dans le mesme temps les circonferences FIKL, GMNO, HPQR, qui sont \uncertain{en mesme raison} que les rayons AF, AG, AH, par lesquelles elles sont descrites. Et pourceque les espaces qui sont parcourus en mesme temps, sont \uncertain{en mesme raison} que les deux vitesses des mobiles, la vitesse \add{\ill{}} \struck{\ill{}} sur toutes les portions du rayon B, F, G, H, seront en mesme raison qu'AB, AF, AG, AH, qui sont aussi les raisons des circonferences.
\note{en marge de gauche, la figure : un cercle BCDE de centre A, coupé par les diamètres CE et BD ; à l'intérieur, en pointillé, les cercles de F, G, H, et, en trait plein, la spirale passant par B, I, N, R jusqu'à A ; lettres C, F, M, P, A, R, Q, L, E sur le diamètre horizontal, G, H, A, K, N, K sur le vertical.}

Donc la vitesse qui est en B diminue \struck{\ill{}} egalement \uncertain{estant egaux} iusques à son repos en A (par le lemme precedent) et la vitesse circulaire totale qui diminue estant reduite à l'uniformité est la moitié de la plus grande vitesse de B. Et si à ce mouvement circulaire \add{de B} il se ioint un autre mouvement \struck{\ill{}} uniforme droit et perpendiculaire vers le centre A, l'helice d'Archimede sera descrite, et le mobile sera porté depuis B, par \struck{\ill{}} le mouvement composé du circulaire \struck{\ill{}} BCDE, et du perpendiculaire BA et viendra à I, d'où il sera porté par le mouvement composé du circulaire FIKL et du droit perpendiculaire \uncertain{CA} iusques à N, et de là iusques à R, par le mouvement composé \struck{\ill{}} du circulaire GMNO, et du droit perpendiculaire DA : et \struck{finalement} du point R il \struck{\ill{}} continuera par un mouvement composé du circulaire HPQR et du droit perpendiculaire EA, iusques au repos en A.

Par cela on void que cette spirale est descrite par le mouvement \struck{circulaire} droit perpendiculaire et par le circulaire qui diminue tousiours iusques au centre A, où il se repose ; c'est à dire par le mouvement circulaire, dont la vitesse reduite à l'uniformité est la moitié de la plus grande vitesse qui est en B, \struck{\ill{}} pendant le temps que le mobile se meut d'un mouvement egal dans le rayon BA.

De plus, soient \struck{la droite} \add{la droite} ba egale à la circonference \add{susdite} BCDE soit iointe à angles droits à la droite ar egale au rayon AB, et soit fait le parallelogramme \uncertain{abrd} divisé par la diagonale br. Cela estant fait, puis qu'un mobile \struck{\ill{}} soit porté \struck{sur} \add{sur} \uncertain{rd} egale au rayon BA du mesme mouvement egal dont nous avons supposé que le mobile se meut sur le rayon AB. Et qu'en mesme temps le mobile se meuve uniformement sur ba dans telle vitesse, qu'en mesme temps que le mobile aura parcouru \uncertain{bd}, il parcourra aussi ba. Donc la vitesse sur ba, qui se fait en mesme temps que le mobile se meut par BA, sera à la vitesse du mobile par la circonference BCDE, comme la droite ba à la circonference BCDE : \struck{et partant} Or ba est la moitié de ladite circonference, et partant la vitesse par ba est la moitié de la vitesse par cette circonference.
\note{les lettres minuscules (ba, ar, br, rd, bd) sont soulignées par l'auteur. En marge de gauche, la figure : un rectangle haut et étroit, lettré en haut b, g, k, n, d, en bas a, c, avec une diagonale et des parallèles lettrées f, e ; l, s ; m, f. Le paragraphe se poursuit, selon toute apparence, au verso (vue 12).}

\page{12}
\note{La vue montre une ouverture : à gauche le verso du feuillet 2 de la vue 11, écrit sur le tiers supérieur seulement, le reste de la page étant blanc ; à droite le recto du feuillet suivant (donné plus bas). En haut de la page de gauche, au milieu, « 8. ». La page de gauche achève la démonstration de la proposition III commencée à la vue 11.}

Mais il a esté prouvé que \add{la vitesse du} \struck{les} mouvemens circulaires qui ont descrit la spirale BINRA, et se diminuent continuellement, est la moitié de la vitesse de la circonference BCDE, donc la vitesse par ba et par la circonference qui se diminue continuellement, et descrivent la spirale, sont egales. Or le mobile est porté \uncertain{depuis} b iusques à e par le mouvement composé des mouvemens par bd et ba, qui se composent \struck{\ill{}} perpendiculairement ; et de e il est porté à s par le mouvement composé des mouvemens par \uncertain{fe}, \uncertain{ge} ; il est porté au point l par le composé \struck{\ill{}} \add{ik, ks} ; et de là il est donc porté à s. Et du point s il est porté par le composé de \uncertain{mf, nf} ; de sorte que la droite br est descrite, laquelle a la puissance de ba et ar, c'est à dire la moitié de la circonference BCDE et le semidiametre BA.
\note{les lettres minuscules sont soulignées par l'auteur ; plusieurs sont écrites vite et se confondent (e et l, s et f, r et e) : les couples de lettres de ce paragraphe sont donnés sous réserve. « ik, ks » est écrit au-dessus d'un groupe biffé.}

Donc puisque la vitesse circulaire se diminuant et la vitesse de BA uniforme sont egales à la vitesse ba et bs, l'une et l'autre à l'autre egale ; et que la spirale et la droite ar sont descrites en mesme temps, (à sçavoir dans le temps que le mobile parcourt BA, ou \uncertain{ar}) et que les angles de la composition sont par tout egaux, et que le composé du droit et du circulaire se trouvant au mesme angle est egal au mouvement composé de ces deux droites, la vitesse sera egale d'une part et d'autre :

Mais les mouvemens \struck{\ill{}} egalement vistes descrivent des lignes egales en temps egaux, donc l'helice BINRA et la droite ar (qui sont descrites par ces mouvemens egaux dans un temps egal) sont egales, et par consequent \struck{la spirale} cette spirale est egale à la droite qui \struck{peut} a la puissance de la circonference et du \struck{\ill{}} semidiametre du cercle.
\note{ici s'arrête le texte de la page ; le reste en est blanc.}

\note{Page de droite de la vue 12 : recto d'un feuillet d'une autre pièce, écrit d'une main plus petite et plus serrée sur toute la page. En haut à droite, à l'encre forte, « 3 » (série continue ; voir l'en-tête) et, plus à droite, d'un trait plus fin, « 10 ». Le cachet rond « BIBLIOTHÈQUE NATIONALE MSS » avec « R.F. » est apposé au milieu de la page. Dans la marge de gauche, en face des deux derniers paragraphes, les chiffres « 1 » et « 2 », et une figure d'ellipse.}

Des triangles Rectilignes et Spheriques. Ch. 1. De la \uncertain{triangulation} \ill{} deux \ill{} \ill{}, 15 pages. Ch. 2. De la cognoiss. des triangles en general, finit à la \uncertain{19} p. Ch. 3. De l'\uncertain{hypothenuse}, finit à la \uncertain{49}. page et suivant \ill{} \ill{} des quarrés iusques \ill{} \uncertain{l'hypothenuse} \ill{} la \uncertain{97} p. C'est \ill{} \ill{} \ill{} et quarrés \uncertain{correspondans}, 2 pages. Ch. 4. De la \ill{} des moindres costés des triangles depuis 1 iusques à la 23 page, où il voit la \uncertain{puissance} qui \ill{} de reste à plusieurs triangles. Ch. 5. Des costés \ill{} des triangles à pag. 23. à pag. 39. Ch. 6. Du costé \uncertain{pair} des triangles. Chap. 7. Des \uncertain{angles} et \uncertain{hypothenuses} des triangles à p. 35 \ill{} ad 63. Chap. 8. De la somme des 2 moindres costés des triangles \ill{} à 91. Chap. 9. De la difference des 2 moindres costés des triangles page 1. Ch. 10 pag. 29. De la somme et de la difference du costé \uncertain{pair} et l'hypothenuse. Chap. 11. p. 51. De la somme et de la difference du costé \uncertain{pair} à l'hypothenuse. Chap. 12. pag. 73. De l'aire des triangles | Chap. 13 à pag. 81. De l'aire des triangles iusques à la page 95. Chap. 14 à pag. 1. Des triangles qui ont \ill{} \ill{} entre les moindres costés. Pag. 9. Chap. 15. Des triangles qui ont \ill{} \ill{} et les costés et qui ont le quarré \ill{} \ill{} entre \ill{} d'un costé \ill{}. P. 27. Chap. 16. Des triangles dont le moindre costé a le quarré \ill{} aux quarrés des 2 autres costés. P. 57. Chap. 17. Des triangles dont la perpendicul. \ill{} sur l'hypot. \ill{} \ill{}. Apres la page 96. il y a des \ill{} qui sont la somme des 3 nombres \uncertain{proportionnels}. Et \ill{} des triangles obliquangles. Chap. 19. Des triangles obliquangles, ausquels \ill{} \ill{} \ill{} \ill{} à p. 3 ad 26. Ch. 20. Des obliquangles qui ont \ill{} \ill{} perpendicul. à pag. 27 iusques à 34. Chap. 21. De \ill{} et \ill{} des triangles obliquangles \uncertain{et p.} 52. De l'aire des obliquangles | Et \ill{} p. 55. Des triangles dont les costés \ill{} \ill{} \ill{} \ill{} \ill{} Ch. 21. Ch. 22. page 57. Des triangles \uncertain{rectangles} \uncertain{isosceles}. Chap. 23. pag. 61. Des triangles equilateres. Chap. \uncertain{24}. Ch. \uncertain{25}. p. \uncertain{65}. Des \uncertain{figures} ou lignes \ill{} qu'on peut \ill{} \struck{dans} \ill{}. De la ligne \ill{} \ill{} dans les Ellipses.
\note{table des matières d'un traité « Des triangles Rectilignes et Spheriques », écrite à la suite sans alinéas, les chapitres séparés par des points et parfois par un trait vertical, que l'on rend par « | ». La main est petite et abrège beaucoup ; les numéros de chapitres et de pages se lisent mieux que les titres, dont de nombreux mots restent illisibles. L'ordre des renvois n'est pas toujours croissant (« page 1 », « à pag. 1 » reviennent au chapitre 9 et au chapitre 14), ce qui suggère plusieurs cahiers paginés séparément ; c'est une inférence. Le chapitre 18 n'est pas nommé à sa place ; « Ch. 21. » est écrit deux fois. Le « 24 » et le « 25 » sont lus avec doute, le 4 et le 7, le 5 et le 6 de cette main se ressemblant. Un trait horizontal sépare cette table de ce qui suit.}

Dans l'Ellipse ACBD. Les principales lignes sont le grand diametre AB, le petit CD, la distance des foyers IL, la ligne IE, et EL \ill{} \ill{} egales au grand diametre \ill{}, de sorte que portant la moitié de AB \ill{} \ill{} \ill{} \add{\uncertain{droit}} C, \ill{} \ill{} \ill{} I et L. \add{IM} \struck{EH}. Et la ligne perpendiculaire au \ill{} aux foyers et EH la perpendiculaire de l'angle E.
\note{en marge de gauche, la figure : une ellipse lettrée A, C, B, D, avec les foyers I et L sur le grand axe, le point E sur la courbe joint à I et à L, les points M, S, K, et des traits pointillés. « IM » est écrit au-dessus de « EH » biffé.}

Pourquoy il y a \ill{} ellipses \ill{}, \uncertain{doutant} que CD peut estre egal à HL, et \ill{} \ill{}, d'autant \ill{} triangles \ill{} fait 2. ellipses qui ont \ill{} \ill{} nombres \uncertain{rationnels} \ill{} par le petit diametre et par la distance des foyers. Comme le grand diametre AB estant 5, le petit pourra estre 4 et la distance des foyers 3 : ou bien le petit diametre pourra estre 3, et IL 4. Lorsque CD et HL sont les 2 costés du triangle dont l'hypothenuse est AB.
\note{un trait horizontal précède ce paragraphe.}

\marginal{1}
Item : trouver une ligne qui \ill{} du grand diametre à tant de \ill{} qu'on voudra et \ill{} \ill{} ellipses rationelles, dont le petit diametre soit \ill{} la distance des foyers \ill{}, et qu'autant le grand que le petit diametre, et la distance des foyers soit rationelle. Par exemple que le grand diametre \ill{} 3 ellipses \ill{}. Il faut adiouster \ill{} \ill{} \ill{} \ill{} \ill{} \ill{} \ill{} \ill{} peut servir à \ill{} \ill{} de triangles rectangles pour donner les 3 ellipses, par exemple, 125, qui se \ill{} à 3 triangles, qui \ill{} pour le grand diametre, et les costés des 3 triangles \struck{\ill{}} \ill{} le petit diametre et la distance des foyers \ill{}. Les 3 triangles sont 44. 117. 125 | 35. 120. 125 | 75. 100 | 125. L'une des ellipses aura 44 pour petit diametre, et 117 pour la distance des foyers ; la \uncertain{seconde} aura 117 pour petit diam. et 44 pour la dist. des foyers et ainsi des 2 autres.
\note{les trois triplets sont exacts : $44^2+117^2 = 35^2+120^2 = 75^2+100^2 = 125^2$. Le troisième est écrit « 75. 100 | 125 », le trait tenant lieu de point.}

\marginal{2}
Une Ellipse estant donnée, avec \ill{} \ill{} nombres rationels, \ill{} demander ses 3 lignes principales, à \uncertain{sçavoir} \ill{} ses grand diametre, \ill{} \ill{} | Le grand costé soit 65, le petit 56, la distance des foyers 33. \uncertain{Trouver} à combien de triangles le grand diametre 65 peut servir d'hypothenuse. \uncertain{Scavoir} à 4 triangles et partant il y a 4 \ill{} \uncertain{ayant} 8 ellipses | les triangles sont \uncertain{16}. 63. 65. 25. 60. 65 | 33. 56. 65. | 39. 52. 65 |
\note{les quatre triplets vérifient $a^2+b^2 = 65^2$ ; le 16 est à demi formé et offert avec doute. Dans « 39. 52. 65 » le 52 est repassé. Un trait horizontal suit.}

La ligne MI estant doublée, est 3\textsuperscript{e} proportionelle au grand et petit diametre, \ill{} CK \ill{} \ill{} \ill{} IM et AK. Donc AK multiplié par MI donne le quarré de CK. Donc le petit et le grand diametre estant donnez, \ill{} MI : car le quarré CK divisé par AK donne MI. \ill{} \ill{} il \ill{} \ill{} \ill{} \add{M}EL \ill{} \ill{} la \ill{} de MI et \ill{} de AB, puisque MI et ML sont egales à AB et partant par la ligne ML \ill{} à MI, \struck{\ill{} EL, \ill{}} : CK. est la R du produit de AK par MI.
\note{« R » est le signe de la racine, un R barré. « MI et ML sont egales à AB » s'entend de leur somme. La page s'arrête ici, sans conclusion ; la vue 13 montre le verso.}

\page{13}
\note{La vue montre une ouverture : à gauche le verso du feuillet de la vue 12 (écrit sur toute la page, de la même main petite et serrée), à droite le recto du feuillet suivant (donné plus bas). En haut de la page de gauche, dans l'angle, un petit signe à l'encre qui n'est pas un chiffre lisible. La numérotation marginale 1, 2 de la vue 12 continue ici : 3, 4, 5, 6.}

AB doit estre le quarré d'un hypotenuse, puisqu'il est quarré \uncertain{et} hypotenuse aussi. Et les 2 lignes MI et ML, doivent \ill{} \ill{} les costés \uncertain{et} \ill{} \ill{} \ill{} des triangles dont \ill{} egalle AB. Il est aisé de \ill{} \ill{} ellipse qui ayt les 4 lignes \uncertain{rationelles} AB, CD, IL, et MI, \ill{} il faut
\note{en marge de gauche, une petite figure : une ellipse lettrée M, C, A, S, K, L, B, D. Plusieurs mots de ces premières lignes sont surchargés.}

Trouver 3 nombres qui \uncertain{soient} des costés aux 2 triangles, \uncertain{Hypotenuse} de \ill{} dispositions soit la somme des \uncertain{hypotenuses} et \ill{} \ill{} costez des triangles. \ill{} \uncertain{nombre} de IL \ill{} \ill{} \ill{} costez au triangle de l'\uncertain{hypothenuse} egal à AB et l'autre costé à CD, \ill{} est aussi de costé à \ill{} triangle IML.

Donc puisqu'il faut que AB soit le quarré d'un hypothenuse, \uncertain{multiplions} \ill{} le triangle \ill{} \ill{} \ill{} \ill{} hypothenuse, comme 3. 4. 5. \ill{} 5. \uncertain{donnera} 15. 20. 25. AB sera donc 25. quarré de 5. Et puisque le quarré de CD doit estre le produit d'AB par le double de \struck{HP} \add{MI} ; il doit estre pair, \uncertain{veu} que le double de \struck{HP} MI doit estre pair, puisqu'il est double du \struck{\ill{}}. Donc \uncertain{puisque} AB \uncertain{est impair}, et le R. de ce produit, \uncertain{à sçavoir} CD, \add{aussi} pair. Le grand des \ill{} 20 sera le petit diametre, et 15 sera la distance des foyers.
\note{« MI » est écrit au-dessus de « HP » biffé ; le second « HP » est aussi biffé. « R. » est le signe de la racine. Un trait horizontal suit ce paragraphe.}

Il faut trouver \uncertain{l'autre} triangle, qui a 15 pour costé et 25 pour la \uncertain{somme} des \ill{} \ill{} \ill{} \ill{}. \uncertain{Partissez} 25 \uncertain{telles} que le quarré \struck{\ill{}} d'une partie \uncertain{soit} \ill{} au quarré de ce \uncertain{premier} \ill{} le quarré de \uncertain{l'autre} partie \ill{} \ill{} au quarré \ill{} \ill{} \ill{} 3. 4. 5. 9. \ill{} le produit de 15 20 25 \ill{} \ill{} \ill{} \ill{} \ill{} les parties donnent 15 et 28. La somme de 3, 4, 5 \ill{} \uncertain{respond} à \uncertain{ces} \ill{} 3 et 5. Donc les \struck{\ill{}} $\square$ sont 9 et 25 | dont la \uncertain{somme} \struck{\ill{}} et la difference sont 34 et 16. Donc les quarrés \ill{} 17 et 8 \uncertain{pour} le costé et \ill{} du triangle \uncertain{2}, qui a 15 pour IL, et les lignes MIL \uncertain{sont} 8. 15. 17. qui a 15 pour costé, et 8 et 15 est aussi de costé au triangle 15. 20. 25. auquel \uncertain{l'hypothenuse} 25 est la somme de 8 et 17. qui sont \ill{} et le costé \ill{} de l'autre triangle. Donc MI sera 8. CD 20. AB 25. IL 15. \struck{\ill{}} Ellipse rationelle \ill{} \ill{} \uncertain{toutes} \uncertain{rationelles} pour MI.
\note{« $\square$ » : un petit carré, signe du quarré. Les triplets 3. 4. 5 ; 15. 20. 25 ; 8. 15. 17 se vérifient ; « 15 et 28 » ne s'accorde pas avec le reste du calcul et est donné tel qu'il se lit.}

\marginal{3.}
Trouver une ligne, à sçavoir pour le grand diametre qui serve à tant d'ellipses qu'on voudra, par exemple à 3, qui ayent les 3 lignes rationelles MI, CD et IL rationelles | prenez le nombre qui est d'hyp. à 3 triangles savoir 125, son quarré est 15625. \uncertain{et} le nombre cherché.
\note{$125^2 = 15625$.}

\marginal{4.}
3 nombres estant donnez desquels à chacun d'ellipse il \ill{}, qui \uncertain{outre} les 3 \ill{} les \ill{} donnez \ill{} \ill{} les 2 autres 4225 et 31250. \uncertain{Lequel} est à 4 ellipses pour \uncertain{parvenir} 65 est d'hyp. à 4 triangles \ill{} \uncertain{c'est} à 6 ellipses, \uncertain{veu} il faut doubler le \uncertain{quotient} des ellipses \ill{} multiplier des 15625 \uncertain{par} 2, lesquels il faut \ill{}
\note{$65^2 = 4225$ ; $2 \times 15625 = 31250$. « 6 ellipses » se lit ainsi, alors que la vue 12 compte 8 ellipses pour 65 : la leçon est donnée telle quelle.}

\marginal{5}
Trouver 3 \uncertain{nombres} grands diametres pour \ill{} \ill{} d'ellipses et plus \uncertain{grande} desquelles ayt la distance des foyers plus grande que le petit diametre et \uncertain{ces} 3 lignes \uncertain{possibles} rationelles. 8450 y \uncertain{satisfait}. \uncertain{Voicy} \ill{} \ill{} \ill{} \ill{} \ill{} \ill{} \ill{} de Carles.
\note{« Carles » est lu tel qu'il se présente ; ce peut être un nom propre. $8450 = 2 \times 65^2$.}

Le \uncertain{nombre} \ill{} \ill{} \ill{} \ill{} \uncertain{quelques} ellip. et le \uncertain{forme} de \uncertain{circulaire}, le \uncertain{autre} et le point \ill{} \ill{} de l'ellipse, qui \uncertain{est} aussi la \ill{} \ill{} \ill{}. Donc
\note{les deux lignes de ce paragraphe sont écrites vite et en partie effacées ; la lecture est très incertaine.}

\marginal{6}
Pour trouver \ill{} \ill{} le grand diam. à tant d'ellipses qu'on voudra en \uncertain{plus} \ill{} \ill{}, \uncertain{à sçavoir} \uncertain{que} \ill{} du \ill{} \ill{}, \uncertain{ausquelles} le petit diametre soit \uncertain{plus grand} que la dist. des foyers, et que le petit diam. de \ill{} et \ill{} du \ill{}, \ill{} et la distance de \ill{} à \uncertain{certain} point \uncertain{soient} \uncertain{rationelles}.
\note{« 6 » est écrit dans la marge de gauche. Au-dessous, deux figures : à gauche une ellipse lettrée A, F, E, C, H, P, O, B, D, I, avec la corde HI perpendiculaire au grand axe ; à droite la même construction en pointillé, lettrée C, F, E, D, H, P, O, B, I, où l'arc HBI est détaché de l'ellipse.}

La \uncertain{surface} HBIO, la \uncertain{portion} \uncertain{quarrée} HBI \uncertain{est} la portion d'ellipse, dont CD est le petit dia. FP la \uncertain{moitié} \ill{} HP la perpendiculaire sur le grand diametre qui de la \uncertain{circonference} HB \ill{} au foyer P. Laquelle ligne HP ou HI \ill{} la distance de l'\uncertain{obiect} FO egal à FH. OB \uncertain{l'espace} du \uncertain{foyer}. \uncertain{Ce} \ill{} :

\marginal{6}
Trouver une ligne \uncertain{rationelle} \uncertain{egale} de grand diam. à tant d'ellipses qu'on voudra en \uncertain{nombre}, \uncertain{à sçavoir} \uncertain{lesquelles} ayent la distance des foyers FP, plus grande que le petit diam. CD, et les lignes CD, FO, OB \uncertain{rationelles}. \uncertain{Car} \ill{} la mesme chose que le probleme precedent. Et puisque FH est quarré et \uncertain{rationel}, OB sera aussi quarré. Lorsque \uncertain{donc} à 4 ellipses \uncertain{servira} le grand dia. \add{\uncertain{de nombres}} 4225 quarré de 65, \uncertain{c'est} 8450 \uncertain{ce}.
\note{le chiffre marginal de ce dernier paragraphe est un « 6 » ouvert qui pourrait être un 7 mal formé ; il est donné comme il se lit. Un signe d'insertion précède « 4225 » et renvoie à deux mots écrits au-dessus, lus avec doute. La page s'arrête ici ; le bas du feuillet est déchiré.}

\note{Page de droite de la vue 13 : recto du feuillet suivant, de la même main que le feuillet 3. En haut à droite, à l'encre forte, « 4 » (série continue ; voir l'en-tête) et, au-dessus, d'un trait fin, « 11 ». La numérotation marginale des problèmes continue : 7, 8.}

Ayant pris le triangle fondamental \uncertain{trouver} toutes les lignes \uncertain{rationelles} de l'ellipse.
\note{en tête de la page, au-dessus d'une figure placée à gauche : une ellipse lettrée A, O, P, E, F, B, C, D, H, avec, en pointillé, un arc de cercle de centre F passant par H et O.}

PCF foyers. HP perpend. sur le foyer. HO portion de cercle dans le rayon est FO ou FH. E centre de l'ellipse.

Que le triangle fondamental soit 8. 15. 17. Le grand diam. 578, sera le double \struck{de l'hyp.} du quarré de l'hypot. 17. Et que FF soit plus grand que CD. FP. se trouve multipliant 17 par le double de 15. | 510. CD multipliant 17 par le double de 8. 272 | PA. multipliant 17 par la differ. de 15 à 17. 34. OA multipliant 15 par la differ. de 15 à 17. | 30 | HF aussi egale FO, estant le quarré de 8 du double du quarré de 17, ou joignant le double du quarré de 15 au quarré de 8, sera 450 à 64. ou ostant 64 de 578 | 514. EA est le quarré de 17. | 289 | PH le quarré de 8. 64 | PO le quarré de la differ. de 15 à 17. 4 | Mais si PF estoit moindre que CD. il faudroit prendre le petit costé 8 au lieu de 15, et 15 au lieu de 8, et partant
\note{« FF » est écrit ainsi ; le sens demande FP. Les calculs sont justes : $2\times17^2 = 578$, $17\times30 = 510$, $17\times16 = 272$, $17\times2 = 34$, $15\times2 = 30$, $2\times15^2+8^2 = 514 = 578-64$. Les traits verticaux « | » séparent les résultats.}

\begin{quote}
PF se trouvera multipliant 17 par le double de 8 --- 272\\
CD. multipl. 17 par le double de 15 --- 510\\
PA multipl. 17 par la differ. de 8 à 17 --- 153\\
OA multipl. 8 par la differ. de 8 à 17 --- 72\\
FO, ou FH ioignant le double du q. de 8 au q. de 15 --- 353\\
EA est q. de 17 --- 289\\
PH q. de 15 --- 225\\
PO le q. de la differ. de 8 à 17 --- 81
\end{quote}
\note{tableau à deux colonnes, les valeurs à droite d'un trait vertical ; les tirets de remplissage de l'auteur sont rendus par des points. Tous les produits sont justes ($17\times9=153$, $8\times9=72$, $2\times64+225=353$).}

PF estant plus grand que CD, le double dont le diametre est FP \ill{} \ill{} \ill{} de l'ellipse et dans son aire peut estre plus grande que celle de l'ellipse, et \ill{} pourra estre plus grand que quelque \uncertain{petit} donné. \ill{} \ill{} \ill{} \ill{} \ill{} \ill{} l'aire \struck{\ill{}} de l'ellipse soit plus grande que \ill{} \ill{}, il faut que le triangle fondamental soit tel, que le produit de l'hypot. par le moindre costé soit plus grand que le q. du grand costé. Donc
\note{paragraphe écrit à droite du tableau, en lignes serrées ; la lecture en est très incomplète.}

\marginal{7}
Trouver 3 triangles, ausquels le produit de l'hypot. par le petit costé soit plus grand que le quarré de l'autre costé.

Ainsi faisant AC est 29. KI. 21. et BD. 20. le \uncertain{cercle} dont le diametre est KI sera KEBI \ill{} \ill{} \ill{} de l'ellipse et aussi aire moindre, \ill{} \ill{} aire \ill{} \ill{} de l'ellipse que 580 à 441. ou ainsi le produit de CA \struck{\ill{}} 29 par BD 20 au q. de KI 21. \ill{} KEI est l'angle droit \ill{} et EN pourra estre rationel et aussi KN. \&c.
\note{en marge de gauche, la figure : une ellipse lettrée C, K, N, L, I, A sur le grand axe, B en haut, F et E sur la courbe, avec un demi-cercle de diamètre KI passant par E, et les droites EK, EI, EN, FK. $29\times20 = 580$, $21^2 = 441$.}

\marginal{8}
Trouver une Ellipse dans laquelle CA, KI, EN, FK. soient rationelles \ill{} \ill{}. La perpend. EN sera rationelle. L'hyp. KI doit estre \uncertain{rationel} q. ou multiplié du q. dont la R. \ill{} hyp. \ill{} le 17 \ill{}. Et si LE egal à LK doit estre rationel, il faut prendre 3 nombres pairs pour l'hyp. IK. \ill{} que \uncertain{la moitié} IK soit rationel. IK sera le double du q. d'une hypoth. Soit par exemple 50. double du 25 q. Le grand costé IE sera 40. et le petit EK 30. et la perpend. EN 24. et puisque AC est \ill{} de IE et EK, il sera 70, prenant la somme de 40 et 30. NK \uncertain{moindre} partie de l'hyp. \ill{} la perpend. EN sera 18. et la grande partie NI 32. | LN. 7. LE. 25. dont le triangle LEN sera 7. \struck{24}. 25. FK \uncertain{sera} 17 1/7. FI 52 6/7. \uncertain{Et} pour oster les fractions, il faut multiplier les lignes \ill{} \ill{} 7. \ill{} aura \ill{} \ill{} toutes les 11 lignes precedentes.
\note{« R. » est le signe de la racine. Le « 24 » de « 7. 24. 25. » est surchargé. $30\cdot40/50 = 24$, $30^2/50 = 18$, $40^2/50 = 32$, $7^2+24^2 = 25^2$.}

\begin{quote}
AC 490\\
IK 350\\
IF 370\\
FK. 120\\
IE 280\\
EK. 210\\
NE. 168\\
IN. 224\\
NK. 126\\
LN 49\\
LE 175.
\end{quote}
\note{colonne de onze valeurs, fermée à droite par un trait vertical : ce sont les onze lignes du problème 8 multipliées par 7. Toutes s'accordent avec le paragraphe précédent ($17\frac17\times7=120$, $52\frac67\times7=370$).}

Il se \uncertain{dit} avec 3 Ellipses la \uncertain{surface} du cercle descrit qui \ill{} perpendiculaire sur le grand diametre. Voy le grand \ill{} à \ill{}.
\note{deux lignes à droite de la colonne de valeurs, au bas de la page ; le reste de la page est blanc.}

\page{14}
\note{La vue montre une ouverture. La page de gauche est le verso du feuillet 4 : blanche, le texte du recto (vue 13) y transparaît à l'envers, le tableau de valeurs à droite. Elle n'est pas transcrite. La page de droite est le recto d'un autre feuillet, d'une main rapide et très raturée, qui paraît être celle des vues 9 à 12 ; le haut du feuillet est bruni et déchiré. En haut à droite, à l'encre forte, « 5 » (série continue ; voir l'en-tête) et, plus à droite, d'un trait fin, « 19 ». Le cachet rond « BIBLIOTHÈQUE NATIONALE MSS » avec « R.F. » est au milieu de la page. Le texte commence au milieu d'un développement : c'est la suite d'une pièce dont le début n'est pas dans ce lot.}

Et nous ne sçaurons pas \add{par reflexion si les rayons de} \struck{\ill{}} lumiere du Soleil quand la terre, ou nos miroirs refleschissent vers le Soleil, vont iusques à la \uncertain{Lune}, quoyque le corps de la Lune privé de la \struck{\ill{}} lumiere directe du Soleil, qui nous la rend claire, monstre \struck{\ill{} \ill{} \ill{}} les \uncertain{inégalités} \uncertain{extremement}, \ill{} \ill{} \ill{} \ill{} \ill{} \ill{} \ill{} \ill{} \ill{} ses corps, \ill{} lesquels les rayons de la terre \ill{} \ill{} \ill{} \uncertain{defaut} si grande force, \struck{\ill{}} \add{qu'ils} nous paroissent fort \struck{\ill{}} obscurément, et quelquefois \ill{} par \uncertain{plusieurs} \ill{} \ill{} \struck{\ill{}} \add{à cause que} la reflexion de nos \uncertain{mers}, \add{et des} autres parties de nostre terre, \struck{\ill{}} n'est pas assez forte pour l'illuminer, \add{\uncertain{n'estoit}} \ill{} \ill{} \ill{} \ill{} \ill{} \ill{} \ill{} \ill{} \ill{} \ill{} \ill{} \ill{} \ill{} \add{\uncertain{sensiblement}} iusques à elle : \struck{\ill{}} du Soleil, qu'elle \struck{\ill{}} \ill{} \ill{}, \struck{\ill{}} peut estre par \uncertain{l'assemblée} de sa reflexion avec \uncertain{une} fois \ill{} \struck{\ill{}} ce qui est difficile à sçavoir, \struck{\ill{}} si l'on s'imagine un oeil qui y fasse l'observation dans \uncertain{sur} la \ill{} \ill{} \ill{} \ill{} \struck{\ill{}} \ill{} \ill{} \ill{} que \struck{\ill{}} \ill{} \ill{} ne doit point que la lumiere que la terre reçoit immediatement du Soleil, \ill{} retourne à ladite Lune, qu'elle illumine grandement.
\note{les quatre premières lignes sont écrites entre des surcharges et des ratures lourdes, avec plusieurs ajouts interlinéaires (« par reflexion si les rayons de », « qu'ils », « à cause que », « des », « sensiblement ») ; l'ordre retenu est une lecture. Dans la marge de gauche, en face de la sixième ligne, « du Soleil » d'une écriture plus appuyée, qui paraît rattachée au texte par un renvoi. Le mot écrit au-dessus de la ligne après « illuminer » se lit « nestoit » ou « n'estoit ».}

S'il ne se perdoit plusieurs rayons, il suivroit aussi de ce que nous voyons que la terre est plus ou moins illuminée, que la Lune \ill{}, soit par la premiere, ou par la \uncertain{seconde} \ill{} \ill{}, ou par la 2 et 3 lumiere du Soleil : mais \struck{\ill{}} \add{\ill{}} inegalités \ill{} \ill{} \ill{} \ill{} \ill{} \ill{} \ill{}, \struck{qu'il \ill{} \ill{} \ill{} \ill{} \ill{}}.
\note{la fin du paragraphe est barrée d'un trait continu sur deux lignes.}

Proposition III.

\add{encor autrement}
Expliquer \uncertain{pourquoy} la reflexion se fait à angles egaux, si l'on \ill{} \ill{} faire la \uncertain{reflexion} perpendiculaire.
\note{« encor autrement » est écrit au-dessus de la ligne avec un signe d'insertion. Une longue ligne horizontale est tirée sous cet énoncé, sur toute la largeur.}

Je \uncertain{pense} que la pensée d'un excellent Philosophe sur \ill{} \struck{\ill{}} \ill{} qu'il \uncertain{desmontra}, c'est à dire \struck{\ill{}} afin qu'il \uncertain{l'estime} \uncertain{embrasse}, et qu'il \uncertain{agréera} davantage, \uncertain{c'est} à dire qu'il \uncertain{justifie} plus raisonnable, et plus veritable.
\note{ligne très surchargée ; « excellent Philosophe » est net, le reste est une lecture douteuse.}

Soit donc une ligne physique AB \struck{\ill{} \ill{} \ill{}}, \struck{\ill{} \ill{} \ill{} \ill{} \ill{} \ill{} \ill{} \ill{}} rigide, qui ne se puisse plier ; et que AB soit \uncertain{meue} \ill{} les 2 perpendiculaires AE, BD, qui seront paralleles et qui aboutiront au \struck{plan reflechissant} \struck{\ill{} \ill{} \ill{} \ill{}} E, D.

Imaginons que cette ligne rigide AB se meuve obliquement, suivant tel angle BDF qu'on voudra, et qu'elle se \uncertain{reflechisse} sur le plan ED \struck{\ill{}} \ill{} \ill{} \ill{} C \uncertain{vers} E \ill{} \ill{} de \uncertain{sorte} \struck{\ill{}} qu'elle \uncertain{sera} droite par \uncertain{demeurant} \add{\uncertain{tenant}} \struck{\ill{}} \ill{} quelque chose de son mouvement la \ill{} ; de sorte que les points C et D \uncertain{laissent} plustost un ovale qu'une ligne droite, quoy que \struck{\ill{}} \ill{} \ill{} \ill{} \ill{} \ill{} ce \uncertain{defaut} de considerer cette particularité, nous \uncertain{retiendrons} le \uncertain{droit}, dans lequel nous supposons que le point D se reflechit iusques à G ; de sorte que le point C arrive aussi tost au \struck{\ill{}} point H du plan reflechissant, que le point D à G.
\marginal{2. M. fol. 86}
\note{la note marginale « 2. M. fol. 86 », d'une encre plus forte et en grands caractères, est encadrée de deux traits horizontaux dans la marge de gauche, en face de ce paragraphe. Elle a l'allure d'un renvoi (un deuxième volume ou manuscrit « M », feuillet 86) ; sa signification n'est pas établie par la pièce. Une ligne horizontale traverse la page au-dessus de « \uncertain{tenant} ».}

Il est evident que le mouvement du point C \add{soit le} \struck{\ill{}} circulaire par HD, ou le perpendiculaire par HI, \ill{} \ill{} par la \uncertain{resistance} du plan reflechissant \uncertain{en} H : Et \uncertain{pourquoy} \ill{} \ill{} que le mouvement circulaire \struck{\ill{}} \ill{} au point H, \uncertain{qu'il} est aussi au point G, qu'il n'y a que le mouvement parallele à la droite ED, \ill{} \ill{} \ill{} \ill{} de la \uncertain{resistance} du plan reflechissant par IH, qui \struck{est} \add{va} perpendiculairement \uncertain{en haut}, qui \uncertain{resiste} au point H ; car il n'y a que le mouvement parallele à ED, et \uncertain{fait} perpendiculaire par KD, qui \uncertain{resiste} au point G. Or \struck{\ill{} \ill{}} le mouvement du point H \uncertain{par} HL, et le mouvement de G par GM, est composé de ces mouvemens.

\struck{\ill{}} Il s'ensuit de ce mouvement composé que \add{\uncertain{l'incidence}} l'angle \struck{\ill{}} \add{BDF} est egal à l'angle \struck{\ill{}} \add{\uncertain{de reflexion} LKE}. \uncertain{Où} il faut remarquer que l'irregularité du mouvement \ill{} du point D vers G, et du point C vers H, \uncertain{vient} de ce que la reflexion \ill{} \ill{} \uncertain{jamais} fait en un point \add{E,} \struck{\ill{}} mais plus prés et au point \add{H :} \ill{}, \struck{\ill{} \ill{} \ill{}} c'est pourquoy cette reflexion se fait par \struck{la droite parallele à \ill{}} \add{une ligne parallele à HI,} semblablement la droite par laquelle se reflechit le point G, est plus proche de C ; autrement GH ne demeureroit pas rigide, et seroit plus longue qu'auparavant.
\note{« BDF » et « LKE » sont écrits au-dessus de groupes de lettres lourdement biffés, précédés de mots interlinéaires (« d'incidence », « de reflexion ») lus avec doute. La page s'arrête ici ; le bas du feuillet est blanc.}

\page{15}
\note{La vue montre une ouverture. À gauche, le verso du feuillet 5, de la même main que la vue 14, dont il continue le texte ; une grosse tache d'encre, en diagonale, couvre le milieu des lignes 3 à 12, et l'écriture du recto transparaît en miroir au bas de la page. À droite, sur un onglet, deux petites figures imprimées, découpées et collées, numérotées à l'encre « 6 » et « 7 », chacune avec le cachet rond de la bibliothèque : la première est un triangle ABC ombré, lettré D, A, I, E, F, G, H, B, C, I, où s'inscrivent des arcs ; la seconde est une figure de courbes symétriques autour d'un axe TD, lettrée P, N, Q, O, T, R, K, X, Y, A, B, E, C, D, L, H, M, I, F, G, avec un cercle et un ovale tangents en A. Elles ne se rapportent pas visiblement au texte de la page et sont décrites ici sans être transcrites. Aucune foliation sur la page de gauche.}

Ce qui a esté demonstré dans la ligne rigide de la longueur d'AB, se demonstre aussi dans la droite \uncertain{semblable} plus grande, et \ill{} d'AB, iusques à l'infini, \ill{} \ill{} \ill{} \ill{} la proposition est vraye dans le point Physique ; de sorte que si AB est un point, les paralleles AE, BD, ne feront qu'une seule ligne droite.

\struck{Dans le \ill{}}
Supposons donc la precedente figure \ill{} \ill{} \ill{} \ill{} \ill{} \ill{} \add{CD} \ill{} \ill{} \ill{} \ill{} CGDH tombe obliquement sur \struck{\ill{}} le plan ED \ill{} les deux paralleles AE, BD : \struck{\ill{}} \ill{} le premier \struck{\ill{}} \ill{} \ill{} \ill{} \ill{} se fera dans un point mathematique, et la premiere impression, dans un point Physique. Posons que \ill{} \ill{} \ill{} continüe, et que la premiere soit la \ill{} \ill{} \ill{} \ill{} \ill{} \ill{} \ill{} le plan HD : le globe \uncertain{s'avancera} \ill{} \ill{} \ill{} les paralleles IH, et KD, \uncertain{selon} la \uncertain{force} de ressort du plan qui se \struck{\ill{}} \add{restablira} dans sa premiere \uncertain{assiette} \ill{}.
\note{la tache d'encre couvre ici trois à cinq mots par ligne, qui sont perdus. « CD » est écrit au-dessus de la ligne.}

Et parce que le globe a \uncertain{receu} le mouvement \uncertain{lateral} vers les parties \ill{} \ill{}, comme auparavant, \uncertain{avec} les paralleles HL, GM, \uncertain{tend} à \struck{\ill{} proportion} \uncertain{pousser} que le mouvement du globe \uncertain{servant} sur le plan DH, vers H, \uncertain{s'establissant} \ill{} \uncertain{restablissant} \uncertain{continuellement} depuis D iusques à H, \uncertain{s'avancera} tout \uncertain{de mesme} : de sorte \struck{\ill{} \ill{} \ill{} \ill{}} qu'il n'y a point d'opposition entre les parties \struck{\ill{}} \uncertain{restablissantes} \struck{\ill{} \ill{} \ill{}} qui se \uncertain{relevent} continuellement devant le globe, qui ne soit recompensée par celles qui s'élevent derriere, \ill{} qui le poussent de mesme force que les parties de devant \struck{\ill{}} s'y opposent.
\marginal{3. M. fol. 87}
\note{la note marginale « 3. M. fol. 87 » est encadrée, dans la marge de gauche, comme « 2. M. fol. 86 » à la vue 14, en face de ce paragraphe. Au milieu du paragraphe, la tache d'encre couvre le mot qui suit « les parties ». « Le restablissement devant le globe ne se fait nullement \ill{} » est une ligne entière barrée avant « qui se relevent ».}

La mesme chose arrive aux \uncertain{petits} \struck{\ill{}} globes dont le mouvement fait la lumiere : et la \struck{\ill{}} raison de la reflexion n'aura plus de difficulté, quand elle se fera obliquement. Quant à la perpendiculaire, il faut imaginer que les \uncertain{corps} tombans et reflechissans, \uncertain{tiennent} par une droite infinie se \uncertain{pressent} en \uncertain{mesme} sorte que deux \ill{} de verre qu'on pousse l'un contre l'autre, qui se \uncertain{rebondissent} et \uncertain{rejaillissent} plus \uncertain{viste} et plus \uncertain{loin} : car \ill{} que les \ill{} \ill{} \ill{} \ill{}, \ill{} la raison \ill{} \ill{} : et parce que le rayon qui illumine, \uncertain{rejaillit} tout autour de soy, et qu'il se fait quelque chose de semblable à \ill{} plustost que tombant sur \ill{} \ill{} \ill{} plate semblable, et \uncertain{rejaillissant}, \struck{\ill{}} \uncertain{fleschie} \ill{} \ill{} à \uncertain{costé}, et en haut ; et que la lumiere qui se reflechit perpendiculairement, fait quelque chose de semblable ; ce qui est \uncertain{aisé} à \uncertain{croire} quand \uncertain{les} \struck{\ill{}} \uncertain{rayons} qui \uncertain{tombent} à plomb, \uncertain{redoublent} leur force : autrement, il est impossible qu'il \uncertain{remonte} tandis qu'il \uncertain{descend}.
\note{paragraphe de lecture très incertaine : la main lie les mots sans lever la plume et en abrège plusieurs ; les mots marqués douteux sont offerts d'après la forme et le sens, et d'autres restent illisibles.}

Avertissement pour la Reflexion des rayons.

Il n'est pas necessaire qu'aucune chose \struck{\ill{}} \add{\uncertain{vivante}} de dessus les \uncertain{miroirs}, pour les \uncertain{reflexions} qu'on y remarque : il suffit que les rayons soient \uncertain{repoussez}, \uncertain{pressez} un \uncertain{moment} \struck{\ill{}} : car la \uncertain{vertu} du \ill{}, qui \uncertain{repousse} \ill{}, est semblable à la reflexion ; \uncertain{bien} qu'il ne \uncertain{rejaillisse} pas comme une bale : \uncertain{pourquoy} la reflexion ne \uncertain{consiste} pas \uncertain{au} \uncertain{rejaillissement}, mais \struck{\ill{}} à la \uncertain{repression}, ou \uncertain{reaction} du corps frappé : de sorte que \struck{\ill{}} \add{tout} \uncertain{rejaillissement} est reflexion ; mais toute reflexion n'est pas un \uncertain{rejaillissement}.

Or cela estant posé, le reste est facile ; parce que les loix de la reflexion sont connües, puisqu'elles sont toutes fondées sur l'égalité des angles : et il n'est pas necessaire de \uncertain{mettre} \ill{} \uncertain{question} si le rayon \uncertain{reflechi} \ill{}, puisque le \ill{} \ill{} \uncertain{repousse} \ill{} \ill{} temps qu'il \struck{\ill{}} \add{\uncertain{rejaillit}} la \uncertain{reflexion} va aussi viste que l'\uncertain{incident} ; puisque \ill{} \ill{} \uncertain{pures}.
\note{la page s'arrête ici ; au-dessous, l'écriture du recto transparaît.}

\page{16}
\note{La vue montre à gauche un onglet et une page blanche, à droite le recto d'un feuillet de la même main rapide qu'aux vues 14 et 15. Les deux petites figures imprimées collées sur onglet, numérotées « 6 » et « 7 » à la vue 15, sont ici rabattues sur le bord gauche de la page et en cachent la marge ; elles ne couvrent pas le texte. En haut à droite, à l'encre, « 8 » (série continue ; voir l'en-tête), et au-dessus, d'un trait fin, « 14 ». Le cachet rond « BIBLIOTHÈQUE NATIONALE MSS » avec « R.F. » est au milieu de la page.}

\uncertain{8} Proposition
\note{le numéro qui précède « Proposition » est un chiffre à longue queue ; il se lit 8, mais la tache qui le couvre en partie laisse le doute.}

Expliquer \struck{\ill{}} toutes les Sections Coniques, à sçavoir la Parabole, l'Ellipse et l'Hyperbole.
\note{titre souligné sur toute la largeur.}

J'ay expliqué ces sections en plusieurs livres, où l'on peut voir leurs proprietez, comme depuis la page 498 du Commentaire sur la Genese iusques à la 535, et dans le 16 chapitre du 4 livre de la verité des sciences ; dans le premier livre de l'Harmonie universelle, prop. 28, et dans la 5 propos. de l'Utilité de l'Harmonie ; et finalement dans la 18 et 19 \add{\uncertain{\&} 17} propos. des Hydrauliques : d'où l'on peut entendre, et \uncertain{prendre} une bonne partie de tout ce qui leur appartient.
\note{renvois lus tels qu'ils se présentent : « 498 », « 535 », « 16 chapitre du 4 livre », « prop. 28 », « 5 propos. », « 18 et 19 ». Au-dessus de « 18 et 19 », un « 17 » précédé d'un signe d'insertion, lu avec doute. Ce paragraphe est écrit à la première personne par l'auteur de ces ouvrages ; la page ne le nomme pas.}

Neantmoins il est necessaire d'en repeter icy quelque chose tant pour l'intelligence de \ill{} \ill{} du \ill{} miroirs, qui sont \uncertain{plus} \ill{} de \uncertain{lentilles} qui suivront. Je \uncertain{mets} donc icy \add{deux} figures qui \add{\uncertain{monstrent}} \struck{\ill{} \ill{}} comme les lignes Coniques qui servent aux miroirs et aux lentilles, naissent \struck{de la \ill{}} \struck{\ill{} \ill{} \ill{} \ill{} particulieres \ill{} \ill{} qui \ill{} \ill{} de la section d'un cone, \ill{} \ill{} \ill{} \ill{}} des differentes coupes ou sections du cone, afin que tous les puissent comprendre sans difficulté par la comparaison de la figure du cone rectangle ABC coupé en trois façons \struck{\ill{}}, qui donnent \struck{\ill{}} l'hyperbole, la parabole et l'ellipse, et de l'autre figure qui est vis à vis, dont A est le point commun par lequel \struck{\ill{}} passent toutes les \ill{} \ill{} \uncertain{sections} \uncertain{separées} de leur cone, \struck{\ill{} \ill{} \ill{}} qui paroissent à pleine veüe, comme les miroirs dont l'axe AD respond droit à l'oeil.
\note{sous « naissent », deux lignes entières sont biffées puis suivies d'une ligne de tirets ; seuls quelques mots s'y lisent. Les « deux figures » annoncées sont, selon toute apparence, les deux figures imprimées numérotées 6 et 7 (vue 15) : le cône coupé ABC et la figure des sections autour du point A et de l'axe TD.}

Je dis donc premierement que la premiere coupe du cone qui se void \struck{\ill{}} \add{dans la figure du cone} GFHE, \struck{GFHE}, \uncertain{et} \ill{} GFHE, dont l'axe \uncertain{est} FE prolongé \uncertain{va} rencontrer \add{au point I} le costé du cone CG prolongé, \struck{\ill{}} et qu'on appelle Hyperbole, est celle qui paroist en HAI dans la figure qui se \struck{\ill{}} void à main droite : ou PRQ \uncertain{son} \uncertain{egale}, \uncertain{estant} comme paralleles \uncertain{accordez}.
\note{« GFHE » est écrit trois fois, deux fois barré ou surchargé. Une grande accolade dans la marge de gauche embrasse ce paragraphe et le suivant.}

Secondement, que la seconde coupe CEHF, \struck{\ill{}} dont l'axe FI est parallele au costé du cone CA, et qu'on appelle \struck{\ill{}} parabole, est FAG dans l'autre figure. En troisiesme lieu, que la coupe FGEH dont l'axe est FE, et qu'on nomme Ellipse, est \struck{\ill{}} l'autre figure, celle \struck{\ill{}} qui a DA pour son grand diametre.

Quant au cercle, qui a B pour son centre dans la 2 figure, il \uncertain{naist} de la coupe \struck{\ill{}} \add{dite} \add{\uncertain{susdite}} qui se fait parallele à la base, \add{\uncertain{comme}} le triangle \uncertain{paroist} de la coupe qui se fait par \uncertain{l'axe} \add{\uncertain{du cone}} ; mais parce qu'ils ont d'autres \uncertain{usages} plus \uncertain{utiles} pour les dioptriques, ie n'en veux point icy parler.
\note{plusieurs mots écrits au-dessus de la ligne, lus avec doute.}

Or cette seconde figure, \uncertain{qui} \struck{\ill{} \ill{}} fait encore voir plusieurs choses \uncertain{remarquables} ; premierement que le point B est le centre du cercle \add{\uncertain{AB}}, \uncertain{et} \uncertain{son} \uncertain{autre} \uncertain{forme} \add{FAG}, le foyer de la parabole, et \struck{\ill{}} l'un des \uncertain{foyers} de l'Ellipse \struck{\ill{}} \uncertain{l'autre} est C : quant à son centre, il est au milieu de la droite tirée de B à C.
\note{interlignes lus avec doute : « AB, et son autre forme », puis « FAG ».}

Le point B est aussi le foyer de l'hyperbole HAI, \uncertain{mais} \ill{} l'appelle \uncertain{interieur} à comparaison de son autre foyer \add{T} qui luy est exterieur. \uncertain{Le} \struck{\ill{}} centre commun de cette hyperbole et de l'autre \uncertain{opposée} PRQ est au point K, auquel se coupent les deux asymptotes MN, LO, \struck{\ill{}} dont la proprieté consiste à ne toucher iamais l'une ou l'autre de ces hyperboles, quelque prolongement qu'on leur puisse donner.
\note{la lettre écrite au-dessus de « foyer » avec un signe d'insertion se lit T. « PRQ » est surchargé d'une tache.}

\struck{L'axe de ces 4 sections}
L'axe de ces 4 sections \struck{\ill{}} est TD : et XY leur \uncertain{sert} \ill{} de tangente, \ill{} qui voudront s'instruire plus particulierement des proprietez de \struck{ces} sections, peuvent lire les 12 qu'\uncertain{ont} \uncertain{mises} dans la 19 proposition de l'Hydraulique : ioint que \struck{\ill{}} ie parleray encore \struck{\ill{} \ill{}} de plusieurs autres propositions.
\note{la page s'arrête ici ; le bas en est blanc.}

\page{17}
\note{La vue 17 et la vue 18 montrent la même ouverture, photographiée deux fois. À gauche, le verso du feuillet 8, écrit de la même main rapide qu'aux vues 14 à 16 ; dans sa marge de gauche sont collées trois petites figures imprimées, découpées, numérotées à l'encre « 9 », « 10 » et « 11 », qui recouvrent ici le début des lignes : « 9 » est une demi-parabole d'axe CBD, les points E reliés à l'axe par des horizontales ; « 10 » un ovale (ellipse) d'axe CBAD, lettré E tout autour, avec des arcs de cercle ; « 11 » la même construction pour une courbe ouverte (hyperbole), lettrée E. À droite, une bande de papier étroite, numérotée « 12 », dépliée hors du volume, porte cinq lignes de la même main. À la vue 18, les figures sont soulevées et le texte est entier : il est transcrit là, et la vue 17 n'est pas transcrite à part.}

\page{18}
\note{Même ouverture que la vue 17, les trois figures collées (9, 10, 11) rabattues hors de la page, le texte découvert. Page de gauche : verso du feuillet 8, aucune foliation. Le texte finit au bas de la page au milieu d'une phrase et se poursuit sur la bande numérotée « 12 », donnée plus bas.}

Proposition.

Expliquer des manieres bien aisées pour descrire les trois sections coniques, à sçavoir la parabole, l'Ellipse et l'Hyperbole.
\note{titre souligné.}

Je ne repete point icy \add{par} quelle coupe, ou section du cone ces sections \uncertain{s'engendrent}, parce que ie les ay assez expliqué dans la prop. precedente, \struck{\ill{} \ill{}} et ceux qui voudront s'instruire plus amplement \uncertain{sur ce subiet}, peuvent lire les sections Coniques du S\textsuperscript{r} Mydorge. Je commence par la parabole, dont ie suppose que le foyer A, et le sommet B soient donnez \struck{\ill{}} ; l'on les peut prendre tels qu'on voudra, suivant les desseins, et la grandeur du miroir qu'on voudra faire.
\note{« S\textsuperscript{r} Mydorge » : le nom est écrit « Mydorge », le titre abrégé avec une lettre suscrite.}

Il faut donc premierement tirer une ligne droite par les points A et B, qui sera \add{\uncertain{ou l'aissieu}} de la parabole, que l'on peut appeller son principal diametre. Et puis il faut faire BC egale à BA. En 3. lieu, il faut prendre sur l'aissieu AB tant de points qu'on voudra ; plus on en prendra, et plus la ligne parabolique sera exacte. En 4 lieu, de chacun de ces points, il faut tracer des perpendiculaires sur ledit aissieu. En 5 lieu, il faut prendre la distance de chaque point D iusques au point C, et le compas estant ouvert de cette distance, il faut mettre l'une de ses iambes au point A, duquel comme centre il faut tracer un arc de cercle ; par le point où cet arc de cercle coupera chaque perpendiculaire qui luy respond, et chaque perpendiculaire, passera la ligne parabolique, comme il arrive dans cette figure qu'elle passe par tous les points E.
\note{les mots écrits au-dessus de « sera » se lisent « ou l'aissieu » ; le mot qu'ils glosent n'a pas été lu. Les trois figures collées dans la marge portent les lettres C, B, D, A, E de cette construction.}

Il faut particulierement remarquer une grande inegalité des points D proche du point B qui est le sommet : et finalement il faut tracer \struck{\ill{}} le plus iustement qu'on peut la parabole depuis le point B d'un costé et d'autre par les points E. Le costé droit de la parabole sera egal au double d'AC, ou au quadruple d'AB ; et tous les rayons du soleil qui tomberont paralleles sur la concavité de cette ligne se reflechiront au point A, qui est le foyer de la parabole, comme \uncertain{nous dirons} \uncertain{apres}, \ill{} expliquant ses proprietez.

L'Ellipse se descrit aussi en supposant ses deux foyers A et F, et le sommet B : qu'on doit prendre suivant l'effet qu'on veut qu'elle produise. Et puis on tire la droite qui passe par A, F, D. \struck{\ill{}} Secondement, soit prise BC egale à BA. En 3. lieu, du point F, comme centre, et de l'intervalle qu'on voudra plus grand que CB, et moindre que FB, soient tracez plusieurs arcs de cercle, comme DE, plus ou moins qu'on voudra, parce que la ligne elliptique approchera plus de la description exacte. En 4. lieu, du centre A, et de l'intervalle de chaque point \struck{\ill{}} D iusques au point C, à sçavoir DC, soit marqué sur chaque arc du cercle qui luy respond, le point d'intersection E, et par tous les points E, et le sommet soit tracée la ligne Elliptique, dont FC est la grandeur de l'aissieu. Nous monstrerons après que tous les rayons qui \ill{} la \struck{\ill{}} ligne \struck{\ill{}} se rassemblent \struck{\ill{}} aux \uncertain{foyers} au point A, ou F, lorsque nous parlerons de ses proprietez.

Il ne reste plus que la troisiesme section du cone à descrire, à sçavoir l'hyperbole, dont s'il faut supposer les deux foyers tels qu'on voudra, par exemple A F. et le sommet B, de sorte que AB soit moindre que la moitié d'AF (car s'il estoit egal, à la moitié, au lieu d'une hyperbole, on traceroit une ligne droite perpendiculaire à AF). Et puis il faut prendre BC egale à BA, et finalement tirer la droite FCBA. En 3. lieu, du \uncertain{centre} F, et de tel intervalle qu'on voudra, plus grand que FB, soient tracez plusieurs arcs de \uncertain{cercle},
\note{le « En 3. lieu » est écrit avec une abréviation raturée. La phrase se poursuit sur la bande numérotée « 12 », à droite de l'ouverture.}

\note{Bande étroite de papier, montée sur onglet et dépliée hors du volume, de la même main. En haut à droite, à l'encre forte, « 12 » (série continue ; voir l'en-tête), et à sa droite, d'un trait fin, « 15 ». Le cachet rond de la bibliothèque est apposé sur la fin des lignes. Elle continue la phrase restée en suspens au bas du verso du feuillet 8.}

comme DE : et du centre A, de l'intervalle de chaque point D iusques au point C, à sçavoir DC, soit marqué sur chaque arc de cercle qui luy respond, le point d'intersection E, et par tous les points E et le sommet soit tracée la ligne hyperbolique, qui servira particulierement pour les refractions, parce qu'elle est propre pour rassembler tous les rayons paralleles dans un mesme point par la refraction qu'elle fait.
\note{un mot est biffé après « par » à la fin de l'avant-dernière ligne ; « refraction » est reporté au début de la ligne suivante.}

\page{20}
\note{La vue 19 n'est pas transcrite : à gauche, le verso blanc de la bande « 12 », le recto y transparaissant ; à droite, une feuille de garde grise, blanche, où transparaissent des cachets. À la vue 20, la page de gauche est le verso de cette feuille grise, blanc ; la page de droite est le recto d'un nouveau feuillet, de la même main rapide, où sont collées deux figures imprimées numérotées à l'encre « 13 » (en haut à gauche : un demi-cercle ACB de centre E, avec le quart de cercle inférieur ADG, lettré L, A, M, I, C, F, H, B, E, N, K, G, D, des pointillés verticaux et les droites FE, HE prolongées) et « 14 » (en bas à gauche : une construction de triangles et de parallèles, lettrée X, N, L, I, C, E, O, D, M, P, Q, K, qui recouvre le début de cinq lignes du texte). En haut à droite, à l'encre forte et soulignée, « 15 » (série continue ; voir l'en-tête), et plus à droite, d'un trait fin, « 16 ». Au-dessus, sur le bord d'un papier qui dépasse la tranche supérieure du feuillet, un « 20 » à l'encre, qui n'appartient pas à ce feuillet. Deux cachets ronds de la bibliothèque sur la page. Le texte de ce feuillet — « 1 Proposition », son corollaire et le début de la « II Proposition » — n'est pas donné ici : le feuillet est photographié de nouveau à la vue 21, béquets relevés, où il se lit en entier ; il est transcrit là, au lot 2.}

\end{document}
